\documentclass[10pt,a4paper]{amsart}
\usepackage{amsmath,amssymb,amsthm,graphicx}
\usepackage{verbatim}
\newenvironment{reference}{\comment}{\endcomment}
%\newenvironment{reference}{}{}
\newenvironment{slogan}{\comment}{\endcomment}
\newenvironment{history}{\comment}{\endcomment}

% For commutative diagrams we use Xy-pic
\usepackage[all]{xy}

% We use 2cell for 2-commutative diagrams.
\xyoption{2cell}
% Theorem environments.
%
\theoremstyle{plain}
\newtheorem{theorem}[subsection]{Theorem}
\newtheorem{proposition}[subsection]{Proposition}
\newtheorem{lemma}[subsection]{Lemma}

\theoremstyle{definition}
\newtheorem{definition}[subsection]{Definition}
\newtheorem{example}[subsection]{Example}
\newtheorem{exercise}[subsection]{Exercise}
\newtheorem{situation}[subsection]{Situation}

\theoremstyle{remark}
\newtheorem{remark}[subsection]{Remark}
\newtheorem{remarks}[subsection]{Remarks}

\numberwithin{equation}{subsection}

% Macros
%
\def\lim{\mathop{\mathrm{lim}}\nolimits}
\def\colim{\mathop{\mathrm{colim}}\nolimits}
\def\Spec{\mathop{\mathrm{Spec}}}
\def\Hom{\mathop{\mathrm{Hom}}\nolimits}
\def\Ext{\mathop{\mathrm{Ext}}\nolimits}
\def\SheafHom{\mathop{\mathcal{H}\!\mathit{om}}\nolimits}
\def\SheafExt{\mathop{\mathcal{E}\!\mathit{xt}}\nolimits}
\def\Sch{\mathit{Sch}}
\def\Mor{\mathop{\mathrm{Mor}}\nolimits}
\def\Ob{\mathop{\mathrm{Ob}}\nolimits}
\def\Sh{\mathop{\mathit{Sh}}\nolimits}
\def\NL{\mathop{N\!L}\nolimits}
\def\CH{\mathop{\mathrm{CH}}\nolimits}
\def\proetale{{pro\text{-}\acute{e}tale}}
\def\etale{{\acute{e}tale}}
\def\QCoh{\mathit{QCoh}}
\def\Ker{\mathop{\mathrm{Ker}}}
\def\Im{\mathop{\mathrm{Im}}}
\def\Coker{\mathop{\mathrm{Coker}}}
\def\Coim{\mathop{\mathrm{Coim}}}

% Boxtimes
%
\DeclareMathSymbol{\boxtimes}{\mathbin}{AMSa}{"02}

%
% Macros for moduli stacks/spaces
%
\def\QCohstack{\mathcal{QC}\!\mathit{oh}}
\def\Cohstack{\mathcal{C}\!\mathit{oh}}
\def\Spacesstack{\mathcal{S}\!\mathit{paces}}
\def\Quotfunctor{\mathrm{Quot}}
\def\Hilbfunctor{\mathrm{Hilb}}
\def\Curvesstack{\mathcal{C}\!\mathit{urves}}
\def\Polarizedstack{\mathcal{P}\!\mathit{olarized}}
\def\Complexesstack{\mathcal{C}\!\mathit{omplexes}}
% \Pic is the operator that assigns to X its picard group, usage \Pic(X)
% \Picardstack_{X/B} denotes the Picard stack of X over B
% \Picardfunctor_{X/B} denotes the Picard functor of X over B
\def\Pic{\mathop{\mathrm{Pic}}\nolimits}
\def\Picardstack{\mathcal{P}\!\mathit{ic}}
\def\Picardfunctor{\mathrm{Pic}}
\def\Deformationcategory{\mathcal{D}\!\mathit{ef}}
\usepackage{cleveref}
\crefname{originalreference}{original reference}{original references}
\Crefname{originalreference}{Original reference}{Original references}
\makeatletter
\newlabel{lemma-point-over-separable-extension}{{0CD4}{source}{}{sourcepointer-1}{}}
\newlabel{lemma-point-over-separable-extension@cref}{{[originalreference][0][]0CD4}{[1][0][]source}{}{}{}}
\makeatother
\crefname{originalreference}{original reference}{original references}
\Crefname{originalreference}{Original reference}{Original references}
\makeatletter
\newlabel{lemma-torsion-picard-smooth-projective}{{0C1Z}{source}{}{sourcepointer-2}{}}
\newlabel{lemma-torsion-picard-smooth-projective@cref}{{[originalreference][0][]0C1Z}{[1][0][]source}{}{}{}}
\makeatother
\crefname{originallemma}{original lemma}{original lemmas}
\Crefname{originallemma}{Original lemma}{Original lemmas}
\makeatletter
\newlabel{pic-lemma-torsion-descends}{{0CD5}{source}{}{sourcepointer-3}{}}
\newlabel{pic-lemma-torsion-descends@cref}{{[originallemma][0][]0CD5}{[1][0][]source}{}{}{}}
\makeatother
\crefname{originallemma}{original lemma}{original lemmas}
\Crefname{originallemma}{Original lemma}{Original lemmas}
\makeatletter
\newlabel{varieties-lemma-smooth-separable-closed-points-dense}{{056U}{source}{}{sourcepointer-4}{}}
\newlabel{varieties-lemma-smooth-separable-closed-points-dense@cref}{{[originallemma][0][]056U}{[1][0][]source}{}{}{}}
\makeatother
\title{The n-torsion of the Picard group of a smooth proper curve becomes the full standard group over a bounded separable extension}
\author{The Stacks Project Authors}
\date{stacks-project:curves.tex:local-snapshot-20261012; GNU FDL 1.2 or later}
\begin{document}
\maketitle
\noindent\textit{Editorial note.} All delivered mathematical text is the author's own, copied from the pinned Stacks Project chapter file curves.tex; no statement or proof step was rewritten, reordered or removed.\par\medskip
\noindent\textit{Editorial note.} The counted proof is the author's complete argument: the bounded rational-point reduction, the two-step torsion descent, the continuous action with open kernel and the factorial Galois bound, and the genus-1 reuse are all retained verbatim.\par\medskip
\noindent\textit{Editorial note.} The label lemma-torsion-picard-becomes-visible names the counted statement, and the proof environment immediately following it in the source is delivered verbatim as its proof.\par\medskip
\noindent\textit{Editorial note.} The counted statement is carried by the author's own lemma environment in the chapter source.\par\medskip
\noindent\textit{Allowed standard prerequisites (editorial).}\begin{itemize}
\item Rational points over bounded separable extensions. Same chapter, Lemma lemma-point-over-separable-extension; Varieties, Lemma varieties-lemma-smooth-separable-closed-points-dense. Supplies a finite separable extension of degree at most 2g - 2 (g >= 2) or merely finite separable (g = 1) over which X acquires a rational point.
\item Torsion over the algebraic closure and its descent. Same chapter, Lemma lemma-torsion-picard-smooth-projective; Picard Schemes of Curves, Lemma pic-lemma-torsion-descends. Identifies Pic(X\_\{overline\{k\}\})[n] with (Z/nZ)\textasciicircum{}\{2g\} and gives the descent of torsion together with the continuous Galois action on Pic(X\_\{k\textasciicircum{}\{sep\}\})[n].
\end{itemize}\medskip
\ifdefined\hypertarget\hypertarget{sourcepointer-1}{}\fi
\noindent\textit{Source reference.} Original reference 0CD4: sources/stacks-curves/pinned-original-source.tex, source lines 2971--2971. Prerequisite pointer invoked at proof line(s) 3660; it names a standard fact of the same chapter that the counted argument uses, and every step of the counted proof built on it is retained verbatim in the delivered body.\par\medskip
\ifdefined\hypertarget\hypertarget{sourcepointer-2}{}\fi
\noindent\textit{Source reference.} Original reference 0C1Z: sources/stacks-curves/pinned-original-source.tex, source lines 3611--3611. Prerequisite pointer invoked at proof line(s) 3666; it names a standard fact of the same chapter that the counted argument uses, and every step of the counted proof built on it is retained verbatim in the delivered body.\par\medskip
\ifdefined\hypertarget\hypertarget{sourcepointer-3}{}\fi
\noindent\textit{Source reference.} Original lemma 0CD5: sources/stacks-chow/tags.txt, source lines 15530--15530. Prerequisite pointer invoked at proof line(s) 3670; it names a standard fact of a sibling chapter that the counted argument uses, and every step of the counted proof built on it is retained verbatim in the delivered body.\par\medskip
\ifdefined\hypertarget\hypertarget{sourcepointer-4}{}\fi
\noindent\textit{Source reference.} Original lemma 056U: sources/stacks-chow/tags.txt, source lines 6722--6722. Prerequisite pointer invoked at proof line(s) 3696; it names a standard fact of a sibling chapter that the counted argument uses, and every step of the counted proof built on it is retained verbatim in the delivered body.\par\medskip
\begin{lemma}
\label{lemma-torsion-picard-becomes-visible}
Let $k$ be a field. Let $n$ be prime to the characteristic of $k$.
Let $X$ be a smooth proper curve over $k$ with $H^0(X, \mathcal{O}_X) = k$
and of genus $g$.
\begin{enumerate}
\item If $g = 1$ then there exists a finite separable extension
$k'/k$ such that $X_{k'}$ has a $k'$-rational point and
$\Pic(X_{k'})[n] \cong (\mathbf{Z}/n\mathbf{Z})^{\oplus 2}$.
\item If $g \geq 2$ then there exists a finite separable extension
$k'/k$ with $[k' : k] \leq (2g - 2)(n^{2g})!$
such that $X_{k'}$ has a $k'$-rational point and
$\Pic(X_{k'})[n] \cong (\mathbf{Z}/n\mathbf{Z})^{\oplus 2g}$.
\end{enumerate}
\end{lemma}

\begin{proof}
Assume $g \geq 2$. First we may choose a finite separable extension
of degree at most $2g - 2$ such that $X$ acquires a rational point, see
Lemma \ref{lemma-point-over-separable-extension}.
Thus we may assume $X$ has a $k$-rational point $x \in X(k)$
but now we have to prove the lemma with
$[k' : k] \leq (n^{2g})!$.
Let $k \subset k^{sep} \subset \overline{k}$ be a separable algebraic
closure inside an algebraic closure.
By Lemma \ref{lemma-torsion-picard-smooth-projective} we have
$$
\Pic(X_{\overline{k}})[n] \cong (\mathbf{Z}/n\mathbf{Z})^{\oplus 2g}
$$
By Picard Schemes of Curves, Lemma \ref{pic-lemma-torsion-descends}
we conclude that
$$
\Pic(X_{k^{sep}})[n] \cong (\mathbf{Z}/n\mathbf{Z})^{\oplus 2g}
$$
By Picard Schemes of Curves, Lemma \ref{pic-lemma-torsion-descends}
there is a continuous action
$$
\text{Gal}(k^{sep}/k)
\longrightarrow
\text{Aut}(\Pic(X_{k^{sep}})[n]
$$
and the lemma is true for the fixed field $k'$ of the kernel of this map.
The kernel is open because the action is continuous which implies
that $k'/k$ is finite. By Galois theory $\text{Gal}(k'/k)$
is the image of the displayed arrow. Since the permutation
group of a set of cardinality $n^{2g}$ has cardinality $(n^{2g})!$
we conclude by Galois theory that $[k' : k] \leq (n^{2g})!$.
(Of course this proves the lemma with the bound
$|\text{GL}_{2g}(\mathbf{Z}/n\mathbf{Z})|$, but all we want
here is that there is some bound.)

\medskip\noindent
If the genus is $1$, then there is no upper bound on the degree of a
finite separable field extension over which $X$ acquires a rational point
(details omitted). Still, there is such an extension for example by
Varieties, Lemma \ref{varieties-lemma-smooth-separable-closed-points-dense}.
The rest of the proof is the same as in the case of $g \geq 2$.
\end{proof}

\end{document}
